\documentclass[11pt]{article}
\usepackage[a4paper,margin=1in]{geometry}
\usepackage[T1]{fontenc}
\usepackage{lmodern}
\usepackage{amsmath,amssymb,amsthm}
\usepackage{booktabs}
\usepackage{array}
\usepackage[hidelinks]{hyperref}
\usepackage{microtype}

\newtheorem{theorem}{Theorem}
\theoremstyle{remark}
\newtheorem{remark}[theorem]{Remark}

\newcommand{\w}[1]{\mathtt{#1}}
\newcommand{\coq}[1]{\texttt{#1}}
\newcommand{\nm}[1]{\mathrm{#1}}

\title{A Transcript Grammar for the Claimed BB(7) Champion\\
\large One move, one gate, and a Coq proof of both --- working notes}
\author{Vishnu Vaidyanathan\\
\small \texttt{vishnu@edgi.in}\\
\small ORCID 0009-0001-6244-5262\\
\small \url{https://github.com/VishnuVaidyanathan/BB5-decimal-macrogrammar}}
\date{11 November 2026}

\begin{document}
\maketitle

\begin{abstract}
These are working notes, not a paper. We study the claimed BB(7) champion
through its transcript, written one digit per step. After a fixed start of
$319$ digits, every digit we can observe is produced by a single move word
$W(P,j)$ applied again and again, and the tape changes by one rule. A second
word describes the halt, and a third a restart. Which of the two happens is
decided by a gate: whether a number $P$ is odd or even. The grammar reproduces
the first $60$ million digits of the run exactly, without running the machine.
We prove in Coq, with no axioms, that the move holds for all values of its
parameters, with its exact step count, and that the gate halts the machine when
$P$ is even. We do not prove the whole run: whether the rule applies until the
end, the halting time, and the claimed bound all remain open.
\end{abstract}

\section{The machine and what is claimed}

\[
\texttt{1RB0RG\_1RC1RB\_0RD1RZ\_0LE0LD\_0LF1LE\_1RA1LF\_1LC0RB}
\]
The machine was found by Jeff Ketchersid in October 2026. The BusyBeaverWiki
lists it as the BB(7) champion with a bound of $3\uparrow^{20}3$ and marks the
claim as not independently verified \cite{wiki}. These notes do not verify that
claim, and the discoverer may already have rules equivalent to the ones below.
Everything here was found from the transcript and from small test tapes.

\paragraph{Digits.} There are $14$ state--symbol pairs. Each step is written as
$2\cdot(\text{state index})+(\text{symbol read})$, with ten to thirteen written
$\w{a},\w{b},\w{c},\w{d}$:

\begin{center}
\begin{tabular}{lccccccc}
\toprule
 & A & B & C & D & E & F & G\\
\midrule
reads 0 & $\w0$ & $\w2$ & $\w4$ & $\w6$ & $\w8$ & $\w{a}$ & $\w{c}$\\
reads 1 & $\w1$ & $\w3$ & $\w5$ & $\w7$ & $\w9$ & $\w{b}$ & $\w{d}$\\
\bottomrule
\end{tabular}
\end{center}

\noindent The digit $\w5$ (C reads 1) is the halt. Three digits repeat in long
runs: $\w3$ (B walks right over ones), $\w9$ and $\w{b}$ (E and F walk left
over ones).

\section{The tape shape}

At certain moments the tape is
\[
1\;0\;[\mathrm{G}\ 0]\;\; 1^{P}\; 0^{g_1}\; 1^{c_1}\; 0^{g_2}\; 1^{c_2}\;\cdots
\]
with the head in state G on a $0$, the two cells $1\,0$ to its left, and
blank cells further left. We call this a \emph{clean shape}. To the right of
the head is a row of blocks of ones separated by gaps of zeros. The first block
has length $P$. A block $1100$ of the earlier notes is simply a block of two
ones followed by a gap of two; we call it a \emph{unit}.

\section{The words}

\begin{center}
\renewcommand{\arraystretch}{1.2}
\begin{tabular}{@{}ll@{}}
\toprule
Name & Digits\\
\midrule
$\nm{Open}$ & $\w{c47698a03}$\\
$\nm{Round}(n)$ & $\w{2476}\,\w9^{n}\,\w{8a0}\,\w3^{n}$\\
$\nm{Build}(P)$ & $\nm{Round}(4)\,\nm{Round}(7)\cdots\nm{Round}(3P+1)$ \quad ($P$ rounds)\\
$\nm{Close}(n)$ & $\w{2468}\,\w{b}^{n}$\\
$\nm{Count}(k)$ & $\w{a1d}\,\w3^{k}\,\w{2468}\,\w{b}^{k+1}$\\
$\nm{Wait}$ & $\w{c476998ba1}$\\
\bottomrule
\end{tabular}
\end{center}

\noindent Every word below starts with the same piece,
\[
\nm{Head}(P,j) \;=\; \nm{Open}\;\nm{Build}(P)\;\nm{Close}(3P+4)\;
\nm{Count}(3P+2)\,\nm{Count}(3P+1)\cdots\nm{Count}(3P+4-2j),
\]
which has $2j-1$ countdown lines. Write $m = 3P+3-2j$.

\paragraph{The move word.}
\[
W(P,j) \;=\; \nm{Head}(P,j)\;\;\w{a1d}\,\w3^{m}\;\nm{Round}(m+1)\;
\w{2476}\,\w9^{m+5}\,\w{8ba1}\;\;\nm{Wait}^{\,j-1}.
\]
Its length is
\[
|W(P,j)| \;=\; 3P^{2}+27P+44+2j+(2j-1)(6P+14-2j).
\]
For $j=1$ this is the five-line block of the earlier notes, of length
$3P^{2}+33P+58$.

\paragraph{The halting word.}
\begin{align*}
H(P,j) \;=\;\; & \nm{Head}(P,j)\;\;\w{a1d}\,\w3^{m}\;\nm{Round}(m+1)\;\nm{Close}(m+5)\\
 & \nm{Count}(m+3)\cdots\nm{Count}(0)\;\;\w{a1c47698b}\;(\w{a1c476998b})^{(m+3)/2}\;\w{ba1d25}.
\end{align*}
It ends in $\w5$, the only $\w5$ of the run.

\paragraph{The restart word} (seen for $j=1$ and $j=2$). It is $H(P,j)$ up to
and including $\w{a1c47698b}$, then $(\w{a1c476998b})^{m/2+1}$, then a fixed
ending that depends only on $j$ (given in the script).

\section{The rule}

A move looks only at $P$, the first gap $g_1=2j$ and the first block $c_1=Z$.

\paragraph{Move} ($Z\ge2$). The machine prints $W(P,j)$ and
\[
1^{P}\;0^{2j}\;1^{Z}\;R \quad\longrightarrow\quad
(1100)^{j-1}\;1^{3P+8-2j}\;00\;1^{Z-2}\;R .
\]
\begin{itemize}
\item \emph{Gap of 2} ($j=1$): $P$ becomes $3P+6$ and the block loses two. This
is the tripling: $P+3$ is multiplied by $3$.
\item \emph{Larger gap}: $j-1$ new units appear in front, so the first block is
now a unit ($P=2$), and the old $P$ becomes a new block of $3P+8-2j$ ones.
\item \emph{When $Z-2=0$} the gap of two that follows joins the next gap, which
therefore grows by two.
\end{itemize}

\paragraph{The gate} ($Z=1$). There is nothing to take two from, and the parity
of $P$ decides.
\begin{center}
\begin{tabular}{@{}l>{\raggedright\arraybackslash}p{0.58\textwidth}@{}}
\toprule
$P$ even & the machine halts and prints $H(P,j)$\\
$P$ odd, and the $1$ is the last cell & restart: it builds $(3P-1)/2$ fresh
units and goes on\\
\bottomrule
\end{tabular}
\end{center}

\section{The run}

The run is
\[
\text{Start}\;\; R(13,2)\;\; W(P_1,j_1)\; W(P_2,j_2)\; W(P_3,j_3)\cdots\; H(P,j).
\]
\begin{itemize}
\item \textbf{Start}: a fixed word of $319$ digits. After it the tape is clean
with $P=13$, a gap of $4$ and a single $1$.
\item \textbf{The gate, first time}: $P=13$ is odd, so the machine restarts and
builds $(3\cdot13-1)/2=19$ units. At step $3{,}954$ the tape is clean with
$P=16$, a gap of $4$, nineteen units and a last single $1$.
\item \textbf{From then on everything is even.} If $P$ is even then $3P+6$,
$3P+8-2j$ and $2$ are even, new units are blocks of two, and blocks only lose
two at a time. So no block can become a single $1$, and the gate is not met
again until the last cell of the tape is reached. At that moment $P$ is even
and the machine halts.
\end{itemize}
So, unlike the mod-3 gate of the BB(5) champion, which is tested at every
layer, this gate is tested once near the start and once at the very end.

\paragraph{The first steps of the real run.}
\begin{center}
\begin{tabular}{rll}
\toprule
Step & Move & Tape after it (blocks and gaps)\\
\midrule
$3{,}954$ & restart $R(13,2)$ & $P=16$;\ gap 4;\ 19 units;\ $1$\\
$5{,}520$ & $W(16,2)$ & $P=2$;\ gap 2, block 52;\ gap 4;\ 18 units;\ $1$\\
$5{,}656$ & $W(2,1)$ & $P=12$;\ gap 2, block 50;\ gap 4;\ 18 units;\ $1$\\
$\approx 6.06\times10^{25}$ & $W(\cdot,1)$, 25 times & $P=15\cdot3^{25}-3$;\ gap 6;\ 18 units;\ $1$\\
\bottomrule
\end{tabular}
\end{center}
The last line is at step $60{,}572{,}642{,}711{,}595{,}382{,}146{,}057{,}206$.

\paragraph{Why it runs so long.} Each time a block is used up, two gaps join
and the next move has a larger $j$. A gap of 2 gives tripling; a gap of 4
repeats a whole run of triplings for each two ones of the next block; and so
on, one level higher for each larger gap. With nineteen units the gaps can grow
to about twenty levels. This fits the shape of the claimed bound, but we have
not worked out or proved any bound.

\section{Tests}

The script \texttt{bb7\_grammar.py} writes the digits from the grammar alone.
\begin{center}
\begin{tabular}{ll}
\toprule
Test & Result\\
\midrule
First $60{,}000{,}000$ digits of the real run & identical to the machine\\
Single moves, $P=2\ldots15$ (odd and even), gaps $2\ldots12$ & 756 of 756 match\\
Complete small runs to the halt & 169 of 169 identical\\
\bottomrule
\end{tabular}
\end{center}
These are tests, not proofs.

\section{What is proved in Coq}

File \texttt{BB7Moves.v} (668 lines, Coq 8.18). \texttt{Print Assumptions}
reports that \coq{move}, \coq{halts\_gen} and \coq{tripling\_chain} are closed
under the global context.

\begin{center}
\renewcommand{\arraystretch}{1.15}
\begin{tabular}{@{}l>{\raggedright\arraybackslash}p{0.72\textwidth}@{}}
\toprule
Name & Statement\\
\midrule
\coq{move} & The move of Section 4, for all $P$, $j\ge1$, $Z\ge2$ and any
rest of the tape, whenever $m=3P+3-2j\ge1$. $P$ may be odd or even.\\
\coq{wcost\_closed} & The move takes exactly $|W(P,j)|$ steps.\\
\coq{triple} & $j=1$: $P\to3P+6$ in $3P^{2}+33P+58$ steps.\\
\coq{chain}, \coq{tri\_closed} & $n$ triplings in a row; $P$ ends as
$(P+3)\,3^{n}-3$.\\
\coq{halts\_gen} & The gate: if $m$ is odd (that is, $P$ even) and the block
after the gap is a single last $1$, the machine halts.\\
\coq{start} & The blank tape reaches the clean shape of step $3{,}954$.\\
\coq{at\_5656} & The tape at step $5{,}656$.\\
\coq{tripling\_chain} & The tape after the 25 triplings, at step
$5656+\coq{chcost 25 12}$.\\
\bottomrule
\end{tabular}
\end{center}

\noindent The proof follows the words: one lemma for each repeated digit
($\w3$, $\w9$, $\w{b}$), one for a round, a countdown line and a waiting line,
and short computations for the fixed pieces. The term \coq{chcost 25 12} equals
the 26-digit step number above by the same arithmetic done outside Coq.

\paragraph{Not proved.}
\begin{enumerate}
\item \textbf{The whole run.} The move needs $m\ge1$, that is, a gap that is
not too large for $P$. We believe the run always satisfies this, because a
small $P$ only appears right after a move, when the gaps in front of it are new
and small; but this is an argument in words. Without it there is no theorem
that the machine halts and no exact halting time.
\item \textbf{The restart} for general odd $P$. The one in the real run is
covered only by running the machine inside Coq up to step $3{,}954$.
\item \textbf{The digit strings.} The step counts equal the lengths of the
words, but the words themselves are not stated as theorems.
\item \textbf{The claimed bound.}
\end{enumerate}

\section{Open questions}

\begin{enumerate}
\item Prove that the run stays inside the rule until the last cell is reached.
This would give a proof of halting and an exact expression for the halting
time, as was done for the BB(6) champion \cite{bb6}.
\item Work out the growth exactly: what is the state when each unit is used
up, and what bound does that give?
\item Prove the restart word for all odd $P$, so that the start is $319$ digits
and nothing else.
\item Explain the $319$-digit start. There is a clean shape already at step
$26$.
\item What does the machine do outside the rule ($m<1$)?
\end{enumerate}

\section*{Acknowledgements}

The machine is due to Jeff Ketchersid. The four readable rules, the repeating
block and the tripling were found by the author by reading the compressed
transcript. The tape analysis, the tests and the Coq development were done with
the help of an AI assistant (Claude, Anthropic); the Coq proofs are checked by
the Coq kernel.

\begin{thebibliography}{9}
\bibitem{wiki}
BusyBeaverWiki. \emph{BB(7)}. \url{https://wiki.bbchallenge.org/wiki/BB(7)},
read 11 November 2026.
\bibitem{bb5}
V. Vaidyanathan. \emph{A Proved Transcript Grammar for the BB(5) Champion}.
November 2026.
\bibitem{bb6}
V. Vaidyanathan. \emph{Exact Step Counts and an Exact Halting Time for the
BB(6) Champion, Proved in Coq}. November 2026.
\end{thebibliography}

\end{document}
